
% --- EJERCICIO 1: PARTÍCULA EN TUBO ROTATORIO ---

\subsubsection{Partícula dentro de un tubo}

Una partícula de masa $m$ se desliza dentro de un tubo como se muestra en la Figura \ref{fig:TuboRotatorio}. El tubo rota en el plano horizontal con velocidad angular constante $\Omega$ en sentido antihorario alrededor del punto fijo $O$. Sabiendo que sobre la partícula actúa una fuerza de fricción viscosa con coeficiente de fricción $c$, y que $r$ es el desplazamiento de la partícula desde el punto $O$.
\begin{enumerate}
    \item determine la ecuación diferencial de movimiento de la partícula en función de $r$.
\end{enumerate}

\begin{figure}[!ht]
    \centering
    \begin{tikzpicture}[>=Latex, scale=1.5]
        
        % --- CONFIGURACIÓN ---
        \def\ang{35}      % Ángulo de inclinación del tubo
        \def\distP{2.4}   % Posición r de la partícula
        \def\Ltubo{3.5}   % Largo total del tubo

        % Base fija O (Suelo pequeño hachurado)
        \fill[pattern=north east lines] (-0.5, -0.2) rectangle (0.5, 0);
        \draw (-0.5, 0) -- (0.5, 0);
        \shadedraw[top color=gris1!60, bottom color=gris1!25, draw=gris1!60] (0.4,0) arc (0:180:0.4) -- cycle;

        % TUBO CON SOMBREADO CILÍNDRICO
        \begin{scope}[rotate=\ang,shift={(0.2,0.2)}]
        
            % Cuerpo del tubo
        % Cuerpo del tubo CON EXTREMO REDONDEADO
            % 1. Relleno con degradado (seguimos el contorno)
            \fill[top color=gris1!50, middle color=gris1!10, bottom color=gris1!50] (0, 0.15) arc (90:270:0.15) -- (\Ltubo, -0.15) -- (\Ltubo, 0.15)-- cycle;

            % 2. Bordes gruesos
            \draw[thick] (0, 0.15) -- (\Ltubo, 0.15);     % Línea superior
            \draw[thick] (0, -0.15) -- (\Ltubo, -0.15);   % Línea inferior
            \draw[thick] (0, 0.15) arc (90:270:0.15); % Tapa redondeada (arco)
            
            % Partícula (Esfera 3D)
            \shadedraw[ball color=black] (\distP, 0) circle (3.5pt) node[above, yshift=8pt] {$m$};
            \fill[black] (0,0) circle (0.06);
            
            % Cota r
   
            \draw[|->] (0, 0.35) -- (\distP, 0.35) node[midway, above] {$r$};
            
            % Etiqueta de Fricción
            \draw[<-] (\distP, -0.2) -- ++(-0.6, -0.6) node[right, font=\small] {Roce viscoso, $c$};
        \end{scope}

        % Punto de pivote O
        \filldraw (0,0) circle (2pt) node[left, xshift=-3pt, yshift=5pt] {$O$};

        % Indicador de velocidad angular Omega
        \draw[->, thick, azul] (0.8, 0.2) arc (0:{\ang}:0.8);
        \node[azul] at (20:1.1) {$\Omega$};
        
        % Referencia horizontal
        \draw[dashed, gris1!60] (0,0.2) -- (1.5, 0.2);

    \end{tikzpicture}
    \caption{Partícula en tubo.}
    \label{fig:TuboRotatorio}
\end{figure}